\documentclass[12pt]{article}
 
\usepackage{fancyhdr}
\usepackage[a4paper, total={6in, 8in}]{geometry}
\usepackage{paralist}
\setlength{\headheight}{15pt} 
\pagestyle{fancy}
\fancyhf{}
\rhead{\today}
\lhead{C. Woodgate}
%\fancyfoot[C]{\thepage}

\begin{document}

\begin{center}
	\large{\textbf{The Idiot's Guide to Recurve Bow Maintenance}}
\end{center}

This is a quick checklist of things every recurve or barebow archer should know
about. Some are measurements you should be able to take/check, while others are
simple maintenance tasks you should be able to perform.

\subsubsection*{Sight Marks}

\begin{itemize}
	\item You will need to set your sight at different heights for different
		distances. This matters most outdoors where we shoot a wide range of
		distances.
	\item The position at which you need to set your sight will also move as
		your technique changes and alterations are made to your bow.
	\item You should {\bf write down} where you set your sight for a given
		distance on a regular basis. This is called a {\bf sight mark}.
\end{itemize}

\subsubsection*{Bracing Height}

\begin{itemize}
	\item This is {\bf the} one thing that every archer should know how to
		check and adjust.
	\item It is the measurement of the perpendicular distance from the bowstring
		to the pressure button.
	\item Adding twists to your string increases it, removing them decreases it.
	\item It should be {\bf the same} every time you shoot your bow; it will
		affect your sight marks and how well your arrows fly.
	\item Measure it, write it down, check it is the same every time you shoot.
	\item As a guide, your BH should fall into one of the following ranges
		depending on your bow length:
		\begin{itemize}
			\item 66'' bow: 210-230mm
			\item 68'' bow: 215-235mm
			\item 70'' bow: 220-240mm
		\end{itemize}
\end{itemize}

\subsubsection*{Waxing Your String}

\begin{itemize}
	\item String waxing serves a few purposes. It helps bond the strands of your
		string together, it helps waterproof your string, and it prolongs the
		life of your string.
	\item You need: 
		\begin{itemize}
			\item String wax
			\item Your string
			\item Your bow
			\item A length of serving thread.
		\end{itemize}
	\item Process:
		\begin{enumerate}
			\item String your bow
			\item Apply a thin layer of wax to string from tube, but not on
				served sections.
			\item Using your fingers, rub the wax into the string. Your fingers
				should get warm but not hot; this helps the wax penetrate the
				fibres.
			\item Wrap the serving thread around the string once, pull it fairly
				tight, and draw it down the string to remove excess wax.
			\item Your string should be nice and round, shiny, and waterproof.
	\end{enumerate}
\end{itemize}

\subsubsection*{Nocking Point}

\begin{itemize}
	\item You should have two nocking points on your string, just over a nock's
		width apart.
	\item These should be {\bf tied on} with serving thread. Brass nocking
		points are bad.
	\item The {\bf top of the bottom nocking point} should sit somewhere between 2 and
		8mm above square, as measured with your bracing height gauge. If you're
		stuck knowing which height to choose, 4mm is a good starting point.
	\item Nocking points sometimes come undone! It is important to know how to
		tie one if this happens to you. There are lots of different ways of
		doing this; learn a way that works for you!
\end{itemize}

\end{document}
